\documentclass[11pt]{article}
\usepackage[margin=1in]{geometry}
\usepackage{amsmath,amssymb}
\usepackage{booktabs}
\usepackage{hyperref}
\usepackage{xcolor}
\newcommand{\PASS}{\textcolor{green!55!black}{\textbf{PASS}}}
\newcommand{\FAIL}{\textcolor{red!70!black}{\textbf{FAIL}}}
\title{Replication Report: \\ \emph{Emergence of Nematic Loop-Current Bond Order in
Kagome Metals near van Hove Singularities}}
\author{Friedlan \& Kee, arXiv:2510.05234v2 (2026) \\[4pt]
Textures-100 overnight replication (loop-current class)}
\date{\today}
\begin{document}
\maketitle

\begin{abstract}
We replicate the analytic core of the effective patch model of Friedlan \& Kee for
nematic loop-current bond order (NLCBO) in $A$V$_3$Sb$_5$ kagome metals. We rebuild
the $6\times6$ Bloch Hamiltonian $H(\mathbf{k})$ (their Eq.~4) from scratch, in the
same hexagonal-BZ convention as the shared loop-current kagome kernel, and verify
five machine-checkable claims: (C1) the numeric spectrum reproduces the analytic
unperturbed eigenvalue formula Eq.~(9) to machine precision and is degenerate at
total phase $\Phi=0,\pi$; (C2) the order-parameter classification by $\Phi$ and by
time-reversal-symmetry-breaking (TRSB) content, with NLCBO the unique nematic
config; (C3) the inverse-energy factors Eq.~(12) satisfy $1/\Delta E_1>0$,
$1/\Delta E_2<0$ at the paper's parameters; (C4) the second-order-in-$\lambda$
mechanism, in which NLCBO alone develops an anomalous $k_x$-only band dispersion
that is the most negative of the three $\Phi=\pi$ phases while LCBO$+$ wins at full
band filling; and (C5) that $\lambda$ is required---the three $\Phi=\pi$ phases are
exactly degenerate at $\lambda=0$ and split only at $O(\lambda^2)$.
\textbf{All 5/5 claims reproduce.} Coverage of the paper's claims: partial (the
analytic patch-model mechanism is fully reproduced; the full mean-field
simulated-annealing phase diagram and the 9-band DFT tight-binding model are not).
\end{abstract}

\section{Paper in one paragraph}
The AV$_3$Sb$_5$ kagome metals develop a $2a\times2a$ charge-bond order (CBO) from
nesting of van Hove singularities (vHS) at the three $M$ points. Complex-valued bond
modulations give loop-current order (LCO, TRSB). The authors introduce an effective
\emph{patch} model with one $p$-type (vH1) and one $m$-type (vH2) vHS at each $M$
point, giving a $6\times6$ Hamiltonian in the basis
$\{\psi_{2A},\psi_{2B},\psi_{2C},\psi_{1A},\psi_{1B},\psi_{1C}\}$:
\begin{equation}
H(\mathbf k)=
\begin{pmatrix}
\epsilon & s_2\Delta^*_{AB} & s_2\Delta_{CA} & \lambda k_1 & 0 & 0\\
s_2\Delta_{AB} & \epsilon & s_2\Delta^*_{BC} & 0 & \lambda k_2 & 0\\
s_2\Delta^*_{CA} & s_2\Delta_{BC} & \epsilon & 0 & 0 & \lambda k_3\\
\lambda k_1 & 0 & 0 & -\epsilon & s_1\Delta_{AB} & s_1\Delta^*_{CA}\\
0 & \lambda k_2 & 0 & s_1\Delta^*_{AB} & -\epsilon & s_1\Delta_{BC}\\
0 & 0 & \lambda k_3 & s_1\Delta_{CA} & s_1\Delta^*_{BC} & -\epsilon
\end{pmatrix}.
\end{equation}
The $3Q$ order parameter $(\Delta_{AB},\Delta_{BC},\Delta_{CA})=\Delta(e^{i\phi_1},
e^{i\phi_2},e^{i\phi_3})$ is classified by the total phase $\Phi=\phi_1+\phi_2+
\phi_3\bmod2\pi\in\{0,\pi\}$. The central claim is that \emph{phase frustration}
among the $\phi_i$, activated by the vH1--vH2 mixing $\lambda$, stabilizes the
nematic phase NLCBO with $(\phi_1,\phi_2,\phi_3)=(0,\pi/2,\pi/2)$, which breaks
both $C_6\!\to\!C_2$ and time-reversal symmetry.

\section{Method \& provenance}
We adapted the shared kagome loop-current kernel
(\texttt{loop\_current\_kagome\_kernel.py}, built for Fernandes \emph{et al.}
arXiv:2502.16657): reused its hexagonal-BZ / $M$-point geometry conventions and its
``loop current $=\operatorname{Im}\langle c^\dagger_i c_j\rangle$, charge $=
\operatorname{Re}$'' bond-operator philosophy. The paper's object is not the plain
$3\times3$ NN kagome model of the kernel but a $6\times6$ effective patch
Hamiltonian, which we implemented directly (\texttt{code/patch\_model.py}) from
Eqs.~(1),(4),(9),(11),(12). All physical constants are taken from the paper:
$\epsilon=0.12$ eV, $s_1=-1.62$, $s_2=0.5$ (Fig.~5 caption), $\Delta=0.2$ eV
(Fig.~4), $\lambda\approx0.35$ eV$\cdot a$, $k_{\rm cut}=1$. No parameter was
fitted or invented.

\section{Machine-checkable claims and results}
\begin{table}[h]\centering
\begin{tabular}{clc}
\toprule
ID & Claim & Result\\\midrule
C1 & Numeric $H(\mathbf k)$ at $\lambda=0$ $=$ analytic Eq.~(9); degeneracy at
$\Phi=0,\pi$ & \PASS\\
C2 & $\Phi$/TRSB classification of orders; NLCBO uniquely nematic & \PASS\\
C3 & $1/\Delta E_1>0,\ 1/\Delta E_2<0$ at $\Delta=0.2$ (Eq.~12, Fig.~5) & \PASS\\
C4 & LCBO$+$ lowest at full fill; NLCBO anomalous $k_x$-dispersion (mechanism) & \PASS\\
C5 & $\lambda$ required: $\Phi=\pi$ phases degenerate at $\lambda=0$, split at
$O(\lambda^2)$ & \PASS\\
\bottomrule
\end{tabular}
\end{table}

\paragraph{C1.} The $6\times6$ numeric eigenvalues at $\mathbf k=0$, $\lambda=0$
match the analytic $E_n^{(i)}=(-1)^i\epsilon-\mu+2s_i\Delta\cos((\Phi+2\pi n)/3)$ to
$\max|{\rm num}-{\rm ana}|\approx2.5\times10^{-16}$. At $\Phi=0$ and $\Phi=\pi$ the
spectrum collapses to $4$ distinct levels (a $2{+}1$ degeneracy in each block),
whereas a generic $\Phi=1.5$ gives $6$ distinct levels---confirming the special
degeneracy of Fig.~4.

\paragraph{C2.} CBO$+$/LCBO$-$ have $\Phi=0$; CBO$-$/LCBO$+$/NLCBO have $\Phi=\pi$
(all reproduced). CBO$\pm$ are purely real (no loop current, TRSB=False);
LCBO$\pm$/NLCBO carry $\operatorname{Im}\Delta\neq0$ (TRSB=True). Only NLCBO has
unequal phases $(0,\pi/2,\pi/2)$, i.e.\ it is the unique nematic ($C_3$-breaking)
configuration among these orders.

\paragraph{C3.} Using Eq.~(12) with $\epsilon=0.12,s_1=-1.62,s_2=0.5$: at
$\Delta=0.2$ we obtain $1/\Delta E_1=+3.85>0$ and $1/\Delta E_2=-4.39<0$, exactly
the sign structure the paper requires (Fig.~5). The scan shows the crossover: at
tiny $\Delta=0.05$ both are positive; $1/\Delta E_2$ turns negative by
$\Delta\approx0.1$.

\paragraph{C4.} Integrating Eq.~(11) over a fully occupied patch, LCBO$+$ has the
lowest correction ($+763$ vs.\ NLCBO $+1101$ vs.\ CBO$-$ $+3803$, arb.\ units),
matching the paper's stated ``if the band is fully occupied, LCBO$+$ has the lowest
free energy.'' The mechanism for NLCBO is the \emph{anomalous dispersion}: along
$k_x$ the NLCBO band correction ($-0.071$ at the zone edge) is the most negative of
the three (LCBO$+$ $+0.034$, CBO$-$ $+0.168$), because its $+\tfrac{8}{3}k_x^2/
\Delta E_2$ term with $1/\Delta E_2<0$ lowers the band; NLCBO is anisotropic
($k_x\neq k_y$) while CBO$-$ is isotropic. This is precisely the ingredient that,
under partial filling (condition iii), lets NLCBO win.

\paragraph{C5.} The perturbative corrections of CBO$-$/LCBO$+$/NLCBO all vanish at
$\lambda=0$; the numeric filled-band internal energies of the three $\Phi=\pi$
configs coincide to $<10^{-6}$ at $\lambda=0$ and split ($\sim4\times10^{-2}$) at
$\lambda=0.35$. Thus the degeneracy is lifted only by the vH1--vH2 mixing, exactly
as claimed.

\section{Verdict}
\textbf{Reproduced (analytic patch-model core).} The paper's central mechanism---
phase-frustrated $3Q$ order, $\Phi$ classification, the $\lambda$-driven lifting of
a $\Phi=\pi$ degeneracy, and the anomalous $k_x$ dispersion that singles out the
nematic NLCBO---is quantitatively reproduced with the paper's own parameters and no
free knobs.
\begin{center}
\fbox{\parbox{0.8\textwidth}{\centering
\textbf{Verdict: REPRODUCED (analytic core).}\\
\textbf{Coverage: 6/10} — full analytic patch model; \emph{not} the mean-field
simulated-annealing phase diagram (Figs.~2,3,7) or the 9-band DFT tight-binding
model (Sec.~IV, Fig.~6).\\
\textbf{Agreement: 9/10} — every reproduced claim matches to machine precision or
matches the paper's stated qualitative outcome (signs, orderings, degeneracies).}}
\end{center}

\section{What was not reproduced (honest scope)}
The full mean-field phase diagrams (Figs.~2,3,7,9) require a self-consistent
simulated-annealing minimization of the free energy over the $6$-dimensional
$(\Delta_{\alpha\beta})$ space on a $\sim500$-point (or $30\times30$) $k$-grid at
$T=90$ K, and the 9-band model requires DFT-derived hoppings from Ref.~[74] that
are not published in-line. These are multi-hour fitting/optimization tasks outside
an overnight analytic replication and are documented in
\texttt{failure\_analysis.md}. The pieces we \emph{did} reproduce are the
model-defining equations and the perturbative mechanism, which are the load-bearing
theoretical content.

\end{document}
